\documentclass[10pt,a4paper]{amsart}
\usepackage[margin=19mm]{geometry}
\usepackage{amsmath,amssymb,mathtools}
\usepackage[scr=rsfs,cal=euler]{mathalfa}
\usepackage{xfrac}
\usepackage[unicode,hidelinks,bookmarks=false]{hyperref}
\newcommand{\ie}{i.e.\ }
\newcommand{\cf}{cf.\ }
\newcommand{\resp}{respectively\ }
\newcommand{\Subsection}{\subsection}
\providecommand{\nobreakdash}{\nobreak\hbox{-}}
\setlength{\emergencystretch}{2em}
\multlinegap0pt
\usepackage{fullpage}
\usepackage{color}
\usepackage{setspace}
\usepackage{here}
\usepackage{float}
\usepackage{mathtools}
\usepackage{latexsym}
\usepackage{tikz-cd}
\newtheorem{Thm}{Theorem}[section]
\newtheorem{Que}[Thm]{Question}
\newtheorem{Lem}[Thm]{Lemma}
\newtheorem{Prop}[Thm]{Proposition}
\newtheorem{Conj}[Thm]{Conjecture}
\newtheorem{Cor}[Thm]{Corollary}
\newtheorem{Cla}[Thm]{Claim}
\newtheorem{Fac}[Thm]{Fact}
\newtheorem{Axi}[Thm]{Axiom}
\newtheorem{Ass}[Thm]{Assumption}
\newtheorem{Def}[Thm]{Definition}
\newtheorem{Rem}[Thm]{Remark}
\newtheorem{Exa}[Thm]{Example}
\def\GL{\mathrm{GL}}
\def\SL{\mathrm{SL}}
\def\SO{\mathrm{SO}}
\def\O{\mathrm{O}}
\def\Sym{\mathrm{Sym}}
\def\sym{\mathrm{sym}}
\def\SPD{\mathrm{SPD}}
\def\dx{{\mathrm{d}x}}
\def\dt{{\mathrm{d}t}}
\def\FR{{\mathrm{FR}}}
\def\Ker{{\mathrm{Ker}}}
\begin{document}
\title[Group invariance of f-divergences and the Fisher--Rao distan]{Group invariance of $f$-divergences and the Fisher--Rao distance}
\author{Frank Nielsen; Sony Computer Science Laboratories Inc.; Kazuki Okamura; Frank Nielsen, Kazuki Okamura}
\date{}
\maketitle
\noindent\emph{Source and licence.} Group invariance of $f$-divergences and the Fisher--Rao distance. Version arXiv:2606.25790v1 (CC BY 4.0). This material is distributed under the Creative Commons Attribution 4.0 International licence, http://creativecommons.org/licenses/by/4.0/, as recorded in the arXiv OAI licence metadata for this exact version and shown on the abstract page https://arxiv.org/abs/2606.25790v1. Full rights notice: licenses/arxiv-2606.25790-CC-BY-4.0.txt.\par\smallskip
\noindent\textit{Editorial scope.} Counted unit: the original \texttt{Thm} environment \texttt{thm:abstract} at source lines 193--200 together with its complete proof at lines 203--210; the delivered document keeps source lines 147--211 so that every referenced label and every citation resolves inside it. Native editable source is the paper's own concatenated TeX; the AMS wrapper is editorial formatting only. No publisher class, upstream script or unrelated artwork is redistributed.
\section{Expression of \texorpdfstring{$f$}{f}-divergences by maximal invariants}

We follow the framework of \cite[Section 3]{Eaton1989}.  
Let $G$ be a topological group. 
Let $(X, \mathcal{B}, \lambda)$ be a $\sigma$-finite measure space. 
Assume that $G$ acts on $X$ measurably, that is, a map $L_g$ on $X$ defined by $L_g (x) \coloneqq   gx$ is measurable. 
Assume that there is a multiplier $\chi$ on $G$ such that 
\begin{equation}\label{eq:rel-inv-mu}
\lambda(gB) = \chi(g) \lambda(B), \quad B \in \mathcal{B}, \quad g \in G. 
\end{equation} 
Let $\Theta$ be a set. 
Assume that $G$ also acts on $\Theta$. 
Let $\{p(\cdot | \theta)\colon\theta \in \Theta\}$ be a family of densities with respect to $\lambda$. 
Let $P_{\theta}(dx) \coloneqq   p(x|\theta)\lambda(dx)$. 
Assume that for every $\theta \in \Theta$, $p(x|\theta) > 0$, $\lambda$-a.e. $x$. 
Assume that 
\begin{equation}\label{eq:inv-density} 
\chi(g) p(gx | g\theta) = p(x|\theta), \ \textup{ $\lambda$-a.e. } x \in X, \quad \theta \in \Theta, \quad g \in G.   
\end{equation}
Let $f$ be a function on $(0,\infty)$ such that $f(1) = 0$ and $f$ is convex on $(0,\infty)$ (and strictly convex at $1$) and 
\[ \int_{X} \left|f \left( \frac{p(x|\theta_2)}{p(x|\theta_1)} \right)\right| p(x|\theta_1) \lambda(dx) < \infty, \quad  \theta_1, \theta_2 \in \Theta. \]
Let the $f$-divergence between $P_{\theta_1}$ and $P_{\theta_2}$ be 
\[ D_f (\theta_1\colon\theta_2) \coloneqq    \int_{X} f \left( \frac{p(x|\theta_2)}{p(x|\theta_1)} \right) p(x|\theta_1) \lambda(dx), \quad  \theta_1, \theta_2 \in \Theta. \]


\begin{Thm}[Group invariance in a transformation model]\label{thm:inv-div}
\[ D_f (g\theta_1 \colon g\theta_2)  = D_f (\theta_1 \colon \theta_2)  \ \  \theta_1, \theta_2 \in \Theta, \ g \in G. \]
\end{Thm}

\begin{proof}
Let 
\[ F(x) \coloneqq   f \left( \frac{p(x|g\theta_2)}{p(x|g\theta_1)} \right) p(x|g\theta_1).  \]
Then, by \eqref{eq:rel-inv-mu}, 
\[ \chi(g) \int_{X} F(gx) \lambda(dx) = \int_{X} F(x) \lambda(dx) =  D_f (g\theta_1 \colon g\theta_2). \]
On the other hand, by \eqref{eq:inv-density}, 
\[ \chi(g) \int_{X} F(gx) \lambda(dx) =  D_f (\theta_1 \colon \theta_2). \]
Thus, we have the assertion. 
\end{proof}

The group $G$ also acts diagonally on the product space $\Theta \times \Theta$ by $(g, (\theta_1, \theta_2)) \mapsto (g\theta_1, g\theta_2)$. 
Let the orbit of $(\theta_1, \theta_2) \in \Theta \times \Theta$ be $O_{(\theta_1, \theta_2)} \coloneqq   \{g(\theta_1, \theta_2) | g \in G\}$. 
If  the action of $G$ on $\Theta \times \Theta$ is transitive, 
then  $O_{(\theta_1, \theta_2)} = \Theta \times \Theta$ and $D_f (\theta_1 \colon \theta_2) = 0$ for every $(\theta_1, \theta_2) \in \Theta \times \Theta$.  
However, even if the action of $G$ on $\Theta$ is transitive,  the action of $G$ on $\Theta \times \Theta$ can be non-transitive. 



\begin{Thm}[Maximal invariant]\label{thm:abstract}
Let $m = m(\theta_1, \theta_2)$ be a self-map on $\Theta \times \Theta$ such that 
$m(\theta_1, \theta_2) \in O_{(\theta_1, \theta_2)}$ for every $(\theta_1, \theta_2) \in \Theta \times \Theta$. 
Assume that $m(g\theta_1, g\theta_2) = m(\theta_1, \theta_2)$ for every $g \in G$ and $\theta_1, \theta_2 \in \Theta$. 
Then the map $m$ is a maximal invariant. 
In particular, 
$D_f (\theta_1 \colon \theta_2) $ is a function of $m(\theta_1, \theta_2)$. 
\end{Thm}

\begin{proof}
If $m(\theta_1, \theta_2) = m(\theta_1^{\prime}, \theta_2^{\prime})$, then 
$O_{(\theta_1, \theta_2)} \cap O_{(\theta_1^{\prime}, \theta_2^{\prime})} \ne \emptyset$. 
Hence, 
$O_{(\theta_1, \theta_2)} = O_{(\theta_1^{\prime}, \theta_2^{\prime})}$. 
Since $(\theta_1^{\prime}, \theta_2^{\prime}) \in  O_{(\theta_1^{\prime}, \theta_2^{\prime})} = O_{(\theta_1, \theta_2)}$, 
there exists $g \in G$ such that $(\theta_1^{\prime}, \theta_2^{\prime}) = g(\theta_1, \theta_2)$. 
\end{proof}


\begin{thebibliography}{99}
\bibitem[Eaton1989]{Eaton1989} Morris~L. Eaton. \newblock {\em Group invariance applications in statistics}, volume~1 of {\em NSF-CBMS Regional Conference Series in Probability and Statistics}. \newblock Institute of Mathematical Statistics, Hayward, CA; American Statistical Association, Alexandria, VA, 1989.
\end{thebibliography}
\end{document}
